%قبل از شروع مسئله ترجیحاً چیزی ننویسید. اما اگر خواستید بنویسید، در قالب لیستی (فقرات و شمارش) بنویسید.
%برای اضافه کردن نمره، بعدی از شروع‌مسئله، [x] را اضافه کنید و به جای x نمره مورد نظر را قرار دهید. برای زیرمسئله‌ها هم به همین صورت است.
%\شروع‌مسئله[10]
\شروع‌مسئله [10]
%\متن‌سیاه{سوال عملی اول:}\\
%\متن‌سیاه{دراپ‌تست}\\

%\متن‌سیاه{ورودی:}\\
%در خط اول ورودی، ابتدا k و سپس m داده می‌شود.\\
%\متن‌سیاه{خروجی:}\\
%در تنها سطر خروجی، تعداد آیفون‌های مورد نیاز می‌آید. اگر در هیچ صورتی چنین کاری ممکن نبود ۱- چاپ می‌شود.\\
%\متن‌سیاه{محدودیت‌ها:}\\
%$1\leq k\leq 10^4$

%$m\leq10^9$
%\\
%\textbf{\متن‌سیاه{ورودی و خروجی نمونه}\\
%% تست اول
%\begin{latin}
% \begin{tabularx}{\textwidth}{|X | X|}
% \hline
% \textbf{stdin} &
% \textbf{stdout}\\
% \hline
% 100 100&
% 1\\
% \hline
% 
% \end{tabularx}
%\end{latin}
%
%% تست دوم
%
%\begin{latin}
% \begin{tabularx}{\textwidth}{|X | X|}
% \hline
% \textbf{stdin} &
% \textbf{stdout}\\
% \hline
% 340 20&
% 3\\
% \hline
% 
% \end{tabularx}
%\end{latin}}
% صورت اصلی مسئله را در این‌جا بنویسید.
مسئله‌ٔ مرتب‌سازی اعداد را در نظر بگیرید. فرض کنید اعداد به صورت دقیق مشخص نیستند. در عوض برای هر عدد، یک بازه‌ٔ بسته روی اعداد حقیقی داریم که می‌دانیم آن عدد به آن بازه متعلق است. یعنی $n$ بازه به صورت $\left[ a_i, b_i \right]$ داریم، که در آن $a_i \le b_i$. 

هدف این است که این بازه‌ها را مرتب کنیم؛ یعنی یک جایگشت $\langle i_1, i_2,..., i_n \rangle$ از بازه‌ها بسازیم به طوری که برای هر\linebreak $j = 1, 2, ..., n$ مقادیر $c_i \in \left[a_{i_j}, b_{i_j}\right]$ موجود باشند که در نامساوی $c_1 \le c_2 \le ...\le c_n$ صدق کنند.

فرض کنید که هر بازه، حداقل با $d-1$ بازه‌ی دیگر همپوشانی داشته باشد. الگوریتمی از $O(n\log \frac{n}{d})$ ارائه دهید که این $n$ بازه با درجه‌ی $d$ را مرتب کند. بدین ترتیب، اگر $d=\Theta(1)$ باشد، الگوریتم در زمان $O(n\log n)$ اجرا می‌شود، و اگر $d=\Theta(n)$ باشد، الگوریتم در زمان $O(n)$ اجرا خواهد شد.
%اگر مسئله چند قسمت داشت به این صورت عمل کنید:
%\شروع{زیرمسائل}
%\زیرمسئله[0] % مسئلۀ ۱آ
%\زیرمسئله % مسئلۀ ۱ب
%\زیرمسئله % مسئلۀ ۱پ
%\زیرمسئله % مسئلۀ ۱ت
%\پایان{زیرمسائل}

\solution{
به این نوع مرتب‌سازی بازه‌ها، مرتب سازی فازی \lr{(Fuzzy Sort)} نیز می‌گویند. الگوریتمی با رویکرد تقسیم و حل و مشابه مرتب‌سازی سریع ارائه می‌دهیم. در هر مرحله از اجرای الگوریتم \lr{Fuzzy Sort}، یکی از بازه‌ها را به عنوان محور (Pivot) قرار می‌دهیم و بازه‌ها را به دو گروه تقسیم می‌کنیم: گروه چپ و گروه راست. الگوریتم \lr{Fuzzy Sort} را به صورت بازگشتی روی این دو دسته صدا می‌زنیم، تا زمانی که تعداد بازه‌های دسته حداکثر $d$ باشد، که در این حالت، بازه‌ها را در زمان $O(d)$ مرتب می‌کنیم.

ناوردای این الگوریتم به این صورت است که در هر مرحله از اجرای الگوریتم روی مجموعه‌ای از بازه‌ها ($S$)، این مجموعه می‌تواند به سه دسته‌ی چپ، وسط، و راست تقسیم کرد، به طوری که شرایط زیر برقرار باشد:
\begin{enumerate}
	\item دسته‌ی چپ، دسته‌ی کوچک‌ترین بازه‌هاست. به بیان دیگر، هر بازه در دسته‌ی چپ، می‌تواند  در ترتیب خروجی، قبل از هر بازه‌ای در دسته‌های وسط و راست قرار بگیرد.
	\item دسته‌ی راست، دسته‌ی بزرگ‌ترین بازه‌هاست. به بیان دیگر، هر بازه در دسته‌ی راست، می‌تواند  در ترتیب خروجی، بعد از هر بازه‌ای در دسته‌های وسط و چپ قرار بگیرد.
	\item در دسته‌ی چپ، هر دو بازه با یک‌دیگر همپوشانی دارند.
	\item در دسته‌ی وسط، هر بازه با حداقل $d-1$ بازه‌ی دیگر (نه لزوماً در دسته‌ی وسط) همپوشانی دارد.
	\item در دسته‌ی راست، هر دو بازه با یک‌دیگر همپوشانی دارند.
	
	
\end{enumerate}
\textbf{شرح الگوریتم:}

همان طور که اشاره شد، این الگوریتم بر پایه‌ی مرتب‌سازی سریع طراحی می‌شود. ابتدا نقاط ابتدایی بازه‌ها (نقاط سمت چپ) را در نظر گرفته و میانه‌ی آن‌ها را در زمان $O(n)$ محاسبه می‌کنیم. بازه‌ی متناظر این نقطه را به عنوان محور در نظر می‌گیریم و بازه‌ها را بر اساس نقاط ابتدایی‌شان partition می‌کنیم (زیرا می‌خواهیم زمان اجرای الگوریتم در بدترین حالت نیز $O(n\log \frac{n}{d})$ باشد). این عمل، مسئله را به دو زیرمسئله تقسیم می‌کند: زیرمسئله‌ی سمت چپ و سمت راست، که اندازه‌ی هر کدام $\lfloor \frac{n}{2} \rfloor$ است.

از آن‌جایی که همه‌ی بازه‌های زیرمسئله‌ی چپ، نقاط ابتدایی کوچک‌‌تری نسبت به محور دارند، همه‌ی آن‌ها می‌توانند در خروجی قبل از محور قرار گیرند. به طور مشابه، تمام بازه‌های زیرمسئله‌ی راست می‌توانند در خروجی بعد از محور قرار گیرند. در نتیجه، \متن‌لاتین{Fuzzy Sort} به صورت بازگشتی روی زیرمسئله‌های چپ و راست به صورت مستقل صدا زده می‌شود.

حال با استفاده از استقرا، ثابت می‌کنیم که ناوردای الگوریتم در هر مرحله از صدا زدن تابع بازگشتی، ثابت می‌ماند. در آغاز، شرایط ناوردا برقرار هستند؛ زیرا می‌توانیم دسته‌ی چپ و راست را خالی بگذاریم و تمام بازه‌ها را در دسته‌ی وسط قرار دهیم.

در ادامه، فرض کنید که ناوردا، در قسمتی از اجرای الگوریتم، برای مجموعه‌ی $S$ از بازه‌ها برقرار باشد. طبق فرض استقرا، می‌توان $S$ را به سه دسته‌ی $S_L$، $S_M$ و $S_R$ افراز کرد (که به ترتیب، همان دسته‌‌های چپ، وسط و راست هستند). برای زیرمسئله سمت چپ محور، سه دسته‌ی ناوردا را مشخص می‌کنیم:

فرض کنید $S^l$ مجموعه‌ی بازه‌های زیرمسئله‌ی چپ باشد. دو حالت پیش می‌آید:
\begin{enumerate}
	\item محور در $S_L$ قرار دارد. قرار می‌دهیم: $S^l_L:=S^L, S^l_R:=\emptyset, S^l_M:=\emptyset$. این سه مجموعه، در شرایط ناوردا صدق می‌کنند.
	\item محور در $S_L$ قرار ندارد. پس می‌توان گفت $S_L \subset S^l$. فرض کنید $T$ مجموعه‌ی بازه‌هایی در $S^l$ باشد که با محور همپوشانی دارند. قرار می‌دهیم: $$S^l_L:=S_L, S^l_R:=T\setminus S^l_L, S^l_M:=S^l\setminus(S^l_L\cup S^l_R)$$ این سه مجموعه، در شرایط ناوردا صدق می‌کنند.
\end{enumerate}
پس ثابت شد که ناوردا برای زیرمسئله‌ی سمت چپ برقرار است. به صورت مشابه، می‌توان برقرار بودن آن را برای زیرمسئله‌ی سمت راست نیز ثابت کرد.

در مرحله‌ی بعد، ادعا می‌کنیم که یک زیرمسئله با تعداد بازه‌های حداکثر $d$ را می‌توان در زمان خطی مرتب کرد. یک مجموعه از بازه‌ها با حداکثر $d$ عضو را در نظر بگیرید. طبق ناوردای الگوریتم، می‌توان بازه‌ها را به سه دسته (که در بالا مشخص شده‌اند) افراز کرد. دو حالت پیش می‌آید. اگر دسته‌ی وسط خالی باشد، یک حالت مرتب‌سازی معتبر این است که که بازه‌های دسته‌ی چپ را به ترتیب دلخواه قرار دهیم، و پس از آن‌ها، بازه‌های دسته‌ی  راست را به ترتیب دلخواه قرار دهیم. و اگر دسته‌ی وسط حداقل یک عضو داشته باشد، این بازه‌ها با تمام $d-1$ بازه‌ی دیگر همپوشانی دارند. پس باز هم می‌توان به سادگی بازه‌ها را مرتب کرد: ابتدا بازه‌هایی که با نقطه‌ی ابتدایی یکی از بازه‌های دسته‌ی وسط $(I)$ همپوشانی دارند، سپس بازه‌ی $I$، و در آخر نیز تمام بازه‌هایی که با نقطه‌ی انتهایی بازه‌ی $I$ همپوشانی دارند را قرار می‌دهیم.

بنابراین ما می‌توانیم درخت بازگشتی را در زمانی که به زیرمسئله‌ای با اندازه‌ی حداکثر $d$ می‌رسیم، هرس کنیم. پس عمق درخت برابر با $\Theta(\lg \frac{n}{d})$ است، و در نتیجه زمان اجرای برنامه $O(n\lg \frac{n}{d})$ است.


}




%%%%%%%%%%%%%%%%%%%%%%%%%%%%%%%%%%%%%%%%%%%%%%%%%%%%%
% برخی دستورات مفید برای تمرین‌ها. %
%%%%%%%%%%%%%%%%%%%%%%%%%%%%%%%%%%%%%%%%%%%%%%%%%%%%%

% فرمول‌های نمونه:
%$f(x) = \Theta(x)$
%$T(x, y) \leq \log(x) + 2^y + \binom{2n}{n}$
% {\tt code\_function}


% لیست موارد بدون شماره‌گذاری:
%\شروع{فقرات}
% \فقره مورد اول 
% \شروع{فقرات}
% \فقره مورد اول داخلی 
% \شروع{فقرات}
% \فقره اولین مورد خیلی داخلی 
% \فقره دومین مورد خیلی داخلی 
% \پایان{فقرات}
% \فقره مورد دوم داخلی 
% \پایان{فقرات}
% \فقره مورد دوم 
%\پایان{فقرات}

% لیست موارد شماره‌گذاری شده:
%\شروع{شمارش}
% \فقره یک مورد 
% \فقره یک مورد دیگر 
% \فقره باز هم یک مورد دیگر
%\پایان{شمارش}

% برای نوشتن یک فرمول چند خطی که یک شماره بخوره. برای حالت‌های دیگر به سایت زیر مراجعه کنید:
% https://www.overleaf.com/learn/latex/Aligning_equations_with_amsmath
%\begin{align} 
% \begin{split}
% (x+y)^3 &= (x+y)^2(x+y) \\
% &= (x^2+2xy+y^2)(x+y) \\
% &= (x^3+2x^2y+xy^2) + (x^2y+2xy^2+y^3) \\
% &= x^3+3x^2y+3xy^2+y^3
% \end{split} 
%\end{align}

% نوشتن ماتریس و آرایه دوبعدی
%\begin{align}
% A = 
% \begin{bmatrix}
% A_{11} & A_{21} \\
% A_{21} & A_{22}
% \end{bmatrix}
%\end{align}

% اضافه کردن تصویر با شرح. برای اضافه کردن تصویر بدون شرح، فقط دستور درج‌تصویر را استفاده کنید.
%\شروع{شکل}[h!]
%\تنظیم‌ازوسط
% \درج‌تصویر[پهنا=0.5\پهنای‌متن]{img.jpg}
%\شرح{شرح تصویر}
%\برچسب{تصویر:۱}
%\پایان{شکل}
%\پایان{زیرمسائل}

%کد زیر را ببینید.
%%\lstinputlisting{source_filename.py}
%\begin{latin}
%\begin{lstlisting}
%# Import the modules
%import sys
%import random
%
%ans = True
%
%while ans:
%question = raw_input("Ask the magic 8 ball a question: (press enter to quit) ")
%
%answers = random.randint(1,8)
%
%if question == "":
%sys.exit()
%
%elif answers == 1:
%print "It is certain"
%
%elif answers == 2:
%print "Outlook good"
%
%elif answers == 3:
%print "You may rely on it"
%
%elif answers == 4:
%print "Ask again later"
%
%elif answers == 5:
%print "Concentrate and ask again"
%
%elif answers == 6:
%print "Reply hazy, try again"
%
%elif answers == 7:
%print "My reply is no"
%
%elif answers == 8:
%print "My sources say no"
%\end{lstlisting}
%\end{latin}